% -------------------------------------------------------------
%    Scrum Implementation
% -------------------------------------------------------------

The project was based on the solid foundation of the Agile methodology, implemented through the Scrum framework. 
Agile Software development follows the general principle of prioritizing people, behaviors, and collaboration within the team members 
rather than following rigid processes, tools and contracts \cite{Adi2015}. 
By its nature is designed to provide a fast, lightweight and alerting environment \cite{Adi2015}. 
In combination with the Scrum framework, a holistic and flexible development product following the Agile principles, 
can provide an increase in speed and flexibility, while enabling iterative development and frequent product increments \cite{Adi2015}.

As this semester’s project requirements were clearly stated, implementing the assignment with the use of the Scrum framework 
would provide both academic and empirical experience. 
Valuable enough to create a solid grounding for each member’s professional competencies; as the knowledge acquired by working 
daily with the scrum framework proved to be educational in many aspects, which will be further analyzed later. 
The team consists of six personnel. 
Despite the number of team members being a rigid requirement from the beginning of the project, 
it is consistent with case studies where small teams are used to maximize performance \cite{Hoegl2005}. 

The team implemented the Scrum Roles as early as the Sprint 1. 
Specifically, the role of the Scrum Master (SM) was decided to be taken in rotation by each member per Sprint. 
All the members were part of the Development Team (DT) which includes programing and reviewing/testing. 
As for the Product Owner (PO), for the first two to three Sprints the role was implemented by the team’s Teacher Assistant, 
as well as the Lecturer in collaboration with the Danfoss representatives. 
As the team members became more familiar with the Product Requirements -for ease of convenience- with a consensus decision 
it was decided to assign the role of PO to one person of the team with the most solid and deep understanding of the requirements. 
Additionally, the team assigned per Sprint, the role of coordinator for the various meetings and ceremonies, which the team held within a week. 

It was decided early on that the team will follow the recommended four to six Sprints in total \cite{Hron2022}. 
It had been agreed upon, that the team would keep the Sprint period at two weeks intervals. 
In the case of an extraordinary circumstance, the Sprint period could be longer, 
but as it has been proved at the end of the project, the group conducted six Sprints in total, 
two weeks each, which concluded in the last week of May.

The industry proposed Sprint ceremonies were reverently followed by the team; Sprint Plannings, 
weekly Stand-up meetings, Sprint Reviews and Sprint Retrospectives \cite{Adi2015}. 
Starting every Sprint with a thorough Sprint Planning. 
nitially, in the first Sprints and while the team was still in the learning process of managing the different Sprint events, 
it was inevitable to stable upon organizational mistakes. 
Thus, due to lack of detailed Sprint Planning in the beginning, incidents like adding tickets mid-Sprint were unavoidable. 
Though, through constructive feedback by the lecturers and intragroup Retrospectives, 
the team managed to create consecutive thorough Sprint Plannings.

The Sprint process continued with weekly Stand-up meetings, which were two at minimum. 
The reason being that the group, while still in the process of learning how to handle the workload of the project, 
were meeting more frequently than normal. 
The actual Stand-up meetings were, as it is being stated in the literature, 
short informative sessions about each other’s work that week \cite{Hoegl2005}. 
The team decided that weekly Stand-up meetings were more efficient for the scale and purpose of this project, rather than daily ones. 
Due to pair-programing, teams of two or three people were meeting more often for the purpose of completing their group’s tickets of the Sprint. 
In other occasions, the team would meet for further organizational purposes, 
reviewing/testing sessions or for solving problems that would need everyone’s attention to be concluded. 
In general, the meetings, as well as the pair-programing and reviewing/testing sessions, were held mostly online through platforms like Discord. 

The Sprint Review and Sprint Retrospectives ceremonies proved to be fruitful for the team. 
In those events many issues, incidents as well as productive ideas were raised. 
As it was expected from a bachelor level project group, it was not possible to avoid behavioral issues, 
schedule conflicts, misaligned work habits and many more. 
Specifically, Sprint 3 had the most frictions between members. 
For that reason, the group implemented the START/STOP/CONTINUE method during the Retrospective, 
contacted with “silent writing” and “round-robin sharing” methods. 
This procedure provided transparency to emotions and grievances and opportunity to open-up, as it offered empowerment and psychological safety. 
Thus, those ceremonies turned out to be tremendously important for the group, 
as through them the team was called upon facing both their individual and collective shortcomings. 
Resulting in crossing the other side of these problems more mature personally and professional than before and the beginning. 

The processes of defining the Product Backlog and Sprint Backlog were forthright and clear-cut. 
The Product Backlog refinement of the requirements can be flagged as the slightly more challenging procedure between the two for the group. 
As the progression of understanding the requirements and solidifying them into concrete 
high level / technical / functional – non-functional requirements were time consuming and complex enough to spent two whole Sprints 
to finalize them and formulate them into UML diagrams, which in result became to be the backbone of the whole code of the project. 
The group did not use the method of translating the requirements into user stories or user cases, 
as it was more efficient based on the team’s working methods dividing the requirements into explicit technical tickets. 

After the clear separation of the work items that needed to be done, the Sprint Backlog procedure for Sprints one until four was straightforward. 
The reason being that those Sprints were clearly technical ones and the work items for those were clearly defined. 
There were also occurrences where the Product Backlog would be emptied and then the team would need to add more tickets 
in the Product Backlog or define more specific ones, but that was not the norm for the group. 
In other cases, incomplete tickets from previous Sprints would be transferred back to the Product Backlog, 
where then the team had a meeting – in the Sprint Planning sessions – deciding which of those would move to the next Sprint 
or would stay there until they become relevant again in a future Sprint.

The chosen tool for implementing the Scrum framework was Jira. 
The reason behind that choice was not a complex one. 
The team considered Jira as one of the more used software for Scrum implementation in the industry, 
thus it was decided to start using it in order to get familiar with it while studying and conquer it before entering the workforce later on. 
For each ticket on Jira, the group used their own developed story points system. 
Specifically, on a 10-point scale, representing the sum of difficulty and importance of a task –both on a scale from one to five. 
The addition of points of difficulty and importance would give the total sum of that individual’s ticket story points. 
Alongside that, each assignee of a ticket had to additionally include the estimated time and the actual time spent on that ticket, 
as well as add the team they were belonging to, in case that ticket needed a pair-programing session.

It was understood that all this information was essential for another tool of Scrum; the burnout charts \cite{Hron2022}. 
The team did not clearly understand in the first two Sprints the importance and usefulness of that tool, 
which can track the productivity of the team by measuring task velocity \cite{Hron2022}. 
If it had been understood, it would be possible to avoid the additional tickets mid-Sprint, 
as by seeing the productivity and progress on the burndown chart, it would help the team plan more efficiently 
future work even from the early stages of production \cite{Hron2022}.

It is evident from the current chapter, deviations from the textbook Scrum framework were inevitable in some cases and purposeful in others. 
Various occurrences of deviation explained beforehand, but it is a common understanding inside the team that those instances were clearly 
justified by two main objectives: ‘Context’ and ‘Performance’ \cite{Hron2022}. 
As it is stated in the source, incidents like the rotation of the Scrum Master role, as well as the implementation of weekly Stand-ups and not daily, 
are being placed into the category ‘Context’, as due to the academic context in which this project is being carried out 
it was unavoidable for schedules and roles to be adapted to the student life. 
Other deviations such as, the customized story point system, as well as holding more weekly meetings except the Scrum ceremonies, 
can be respectively categorized into the ‘Performance’ category, where limited time made the team deviate from the traditional path.

As the sources mention, following all the Scrum events influences the perception of being successful at Scrum \cite{Kadenic2023}. 
To complement on that idea, it is noteworthy that for the group following all the Scrum events had more of a sense of being successful 
as an individual working towards being a more mature professional and a person; learning by personal mistakes, adapting to challenges, 
and supporting continuous learning and improvement \cite{Hron2022}.
